\section{AMC 12 12B 2012}
        \begin{enumerate}
        	\item (\textit{AMC 12 12B 2012}) Each third-grade classroom at Pearl Creek Elementary has 18 students and 2 pet rabbits. How many more students than rabbits are there in all 4 of the third-grade classrooms?


 $\textbf{(A)}\ 48\qquad\textbf{(B)}\ 56\qquad\textbf{(C)}\ 64\qquad\textbf{(D)}\ 72\qquad\textbf{(E)}\ 80$

 \textbf{Answer: }C



\textbf{Solution 1}

Multiplying $18$ and $2$ by $4$ we get $72$ students and $8$ rabbits. We then subtract: $72 - 8 = \boxed{\textbf{(C)}\ 64}.$




\textbf{Solution 2}

In each class, there are $18-2=16$ more students than rabbits. So for all classrooms, the difference between students and rabbits is $16 \times 4 = \boxed{\textbf{(C)}\ 64}$



	\item (\textit{AMC 12 12B 2012}) A circle of radius 5 is inscribed in a rectangle as shown. The ratio of the length of the rectangle to its width is 2:1. What is the area of the rectangle?

                    \begin{figure}[H]
                        \centering
                        \includegraphics[width=0.5\textwidth]{Problem/AMC_12_12B_2012_2_Problem_Figure_1.png}
                    \end{figure}
                    
$\textbf{(A)}\ 50\qquad\textbf{(B)}\ 100\qquad\textbf{(C)}\ 125\qquad\textbf{(D)}\ 150\qquad\textbf{(E)}\ 200$

 \textbf{Answer: }E



\textbf{Solution}

If the radius is $5$, then the width is $10$, hence the length is $20$. $10\times20= \boxed{\textbf{(E)}\ 200}.$



	\item (\textit{AMC 12 12B 2012}) For a science project, Sammy observed a chipmunk and squirrel stashing acorns in holes. The chipmunk hid 3 acorns in each of the holes it dug. The squirrel hid 4 acorns in each of the holes it dug. They each hid the same number of acorns, although the squirrel needed 4 fewer holes. How many acorns did the chipmunk hide? 


 $\textbf{(A)}\ 30\qquad\textbf{(B)}\ 36\qquad\textbf{(C)}\ 42\qquad\textbf{(D)}\ 48\qquad\textbf{(E)}\ 54$

 \textbf{Answer: }D



\textbf{Solution}

If $x$ is the number of holes that the chipmunk dug, then $3x=4(x-4)$, so $3x=4x-16$, and $x=16$. The number of acorns hidden by the chipmunk is equal to $3x = \boxed{\textbf{(D)}\ 48}$




\textbf{Solution 2}

Trying answer choices, we see that $\boxed{\textbf{(D)}\ 48}$ works. ~Extremelysupercooldude

 


	\item (\textit{AMC 12 12B 2012}) Suppose that the euro is worth $1.30$ dollars. If Diana has $500$ dollars and Etienne has $400$ euros, by what percent is the value of Etienne's money greater than the value of Diana's money?


 $\textbf{(A)}\ 2\qquad\textbf{(B)}\ 4\qquad\textbf{(C)}\ 6.5\qquad\textbf{(D)}\ 8\qquad\textbf{(E)}\ 13$

 \textbf{Answer: }B



\textbf{Solution}

The ratio $\frac{400 \text{ euros}}{500 \text{ dollars}}$ can be simplified using conversion factors:
 \[\frac{400 \text{ euros}}{500 \text{ dollars}} \cdot \frac{1.3 \text{ dollars}}{1 \text{ euro}} = \frac{520}{500} = 1.04\] which means the money is greater by $\boxed{ \textbf{(B)} \ 4 }$ percent.

 



\textbf{Solution 2}

If we divide each of Etienne's and Diana's values by $100$, the problem stays the same. Then, Etienne has $1.3$ times the amount of money Diana has, so Etienne has $5.2$ dollars. Since $\dfrac{5.2}{5} = 1.04$, Etienne has $\boxed{ \textbf{(B)} \ 4 }$ percent more money than Diana. ~Extremelysupercooldude

 


	\item (\textit{AMC 12 12B 2012}) Two integers have a sum of 26. When two more integers are added to the first two, the sum is 41. Finally, when two more integers are added to the sum of the previous 4 integers, the sum is 57. What is the minimum number of even integers among the 6 integers? 


 $\textbf{(A)}\ 1\qquad\textbf{(B)}\ 2\qquad\textbf{(C)}\ 3\qquad\textbf{(D)}\ 4\qquad\textbf{(E)}\ 5$

 \textbf{Answer: }A



\textbf{Solution}

Since, $x + y = 26$, $x$ can equal $15$, and $y$ can equal $11$, so no even integers are required to make 26. To get to $41$, we have to add $41 - 26 = 15$. If $a+b=15$, at least one of $a$ and $b$ must be even because two odd numbers sum to an even number. Therefore, one even integer is required when transitioning from $26$ to $41$. Finally, we have the last transition is $57-41=16$. If $m+n=16$, $m$ and $n$ can both be odd because two odd numbers sum to an even number, meaning only $1$ even integer is required. The answer is $\boxed{\textbf{(A)}}$. ~Extremelysupercooldude (Latex, grammar, and solution edits)

 



\textbf{Solution 2}

Just worded and formatted a little differently than above.

 The first two integers sum up to $26$. Since $26$ is even, in order to minimize the number of even integers, we make both of the first two odd.

 The second two integers sum up to $41-26=15$. Since $15$ is odd, we must have at least one even integer in these next two.

 Finally, $57-41=16$, and once again, $16$ is an even number so both of these integers can be odd.

 Therefore, we have a total of one even integer and our answer is $\boxed{\textbf{(A)}}$.

 


	\item (\textit{AMC 12 12B 2012}) In order to estimate the value of $x-y$ where $x$ and $y$ are real numbers with $x > y > 0$, Xiaoli rounded $x$ up by a small amount, rounded $y$ down by the same amount, and then subtracted her rounded values. Which of the following statements is necessarily correct?


 $\textbf{(A)}\ \text{Her estimate is larger than }x-y$
$\textbf{(B)}\ \text{Her estimate is smaller than }x-y$
$\textbf{(C)}\ \text{Her estimate equals }x-y$
$\textbf{(D)}\ \text{Her estimate equals }y - x$
$\textbf{(E)}\ \text{Her estimate is 0}$

 \textbf{Answer: }A



\textbf{Solution}

The original expression $x-y$ now becomes $(x+k) - (y-k)=(x-y)+2k>x-y$, where $k$ is a positive constant, hence the answer is $\boxed{\textbf{(A)}}$.

 



\textbf{Solution 2}

The problem never says what $x$ and $y$ are, so we can decide what they are. Let $x = 1.6$ and $y = 1.4$. We round $x$ to $2$ and $y$ to $1$. Then the new $x - y = 1$, while the original $x - y = 0.2$. Thus, the new $x - y$ is greater than the original $x - y$. The answer is $\boxed{\textbf{(A)}}$. ~Extremelysupercooldude

 


	\item (\textit{AMC 12 12B 2012}) Small lights are hung on a string 6 inches apart in the order red, red, green, green, green, red, red, green, green, green, and so on continuing this pattern of 2 red lights followed by 3 green lights. How many feet separate the 3rd red light and the 21st red light?


 \textbf{Note:} 1 foot is equal to 12 inches.


 $\textbf{(A)}\ 18\qquad\textbf{(B)}\ 18.5\qquad\textbf{(C)}\ 20\qquad\textbf{(D)}\ 20.5\qquad\textbf{(E)}\ 22.5$

 \textbf{Answer: }E



\textbf{Solution}

We know the repeating section is made of $2$ red lights and $3$ green lights. The 3rd red light would appear in the 2nd section of this pattern, and the 21st red light would appear in the 11th section. There would then be a total of $44$ lights in between the 3rd and 21st red light, translating to $45$ $6$-inch gaps. Since the question asks for the answer in feet, the answer is $\frac{45*6}{12} \rightarrow \boxed{\textbf{(E)}\ 22.5}$.

 



\textbf{Solution 2}

We know that both $3$ and $21$ are odd. This means that they both start their respective patterns of $rrggg$ as the first $r$ value.

 Let's take a look at one full gap of two reds $rrgggr$, in this gap there is a total of $2.5$ feet of gap (considering that $6$ inches is half a foot). So for each gap of two reds there is a $2.5$ feet gap.

 The total amount of red gaps of lights hung is $\frac{21-3}{2} = 9$ gaps. Multiplying $9 * 2.5 = 22.5$ inches. This gives $(E)$.

 


	\item (\textit{AMC 12 12B 2012}) A dessert chef prepares the dessert for every day of a week starting with Sunday. The dessert each day is either cake, pie, ice cream, or pudding. The same dessert may not be served two days in a row. There must be cake on Friday because of a birthday. How many different dessert menus for the week are possible?


 $\textbf{(A)}\ 729\qquad\textbf{(B)}\ 972\qquad\textbf{(C)}\ 1024\qquad\textbf{(D)}\ 2187\qquad\textbf{(E)}\ 2304$

 \textbf{Answer: }A



\textbf{Solution}

We can count the number of possible foods for each day and then multiply to enumerate the number of combinations.

 On Friday, we have one possibility: cake.

 On Saturday, we have three possibilities: pie, ice cream, or pudding. This is the end of the week.

 On Thursday, we have three possibilities: pie, ice cream, or pudding. We can't have cake because we have to have cake the following day, which is the Friday with the birthday party.

 On Wednesday, we have three possibilities: cake, plus the two things that were not eaten on Thursday.

 Similarly, on Tuesday, we have three possibilities: the three things that were not eaten on Wednesday.

 Likewise on Monday: three possibilities, the three things that were not eaten on Tuesday.

 On Sunday, it is tempting to think there are four possibilities, but remember that cake must be served on Friday. This serves to limit the number of foods we can eat on Sunday, with the result being that there are three possibilities: The three things that were not eaten on Monday.

 So the number of menus is $3 \cdot 3 \cdot 3 \cdot 3 \cdot 3 \cdot 1 \cdot 3 = 729.$The answer is $\boxed{A}$.

 



\textbf{Solution 2}

We can perform casework as an understandable means of getting the answer. We can organize our counting based on the food that was served on Wednesday, because whether cake is or is not served on Wednesday seems to significantly affect the number of ways the chef can make said foods for that week.

 Case 1: Cake is served on Wednesday. Here, we have three choices for food on Thursday and Saturday since cake must be served on Friday, and none of these choices are cake, which was served Wednesday. Likewise, we have three choices (pie, ice cream, and pudding) for the food served on Tuesday and thus three choices for those served on Monday and Sunday, with these three choices being whatever was not served on Tuesday and Monday, respectively. Hence, for this case, there are $3 \cdot 3 \cdot 3 \cdot 3 \cdot 3 = 243$ possibilities.

 Case 2: Cake is not served on Wednesday. Obviously, this means that pie, ice cream, and pudding are our only choices for Wednesday's food. Since cake must be served on Friday, only ice cream, pudding, and cake can be served on Thursday. However, since one of those foods was chosen for Wednesday, we only have two possibilities for Thursday's food. Like our first case, we have three possibilities for the food served on Tuesday, Monday, and Sunday: whatever was not served on Wednesday, Tuesday, and Monday, respectively. $3 \cdot 3 \cdot 3 \cdot 3 \cdot 2 \cdot 3 = 486$ possibilities thus exist for this case.

 Adding the number of possibilities together yields that $243 + 486 = 729$ is the total number of menus, making our answer $\boxed{A}$.

 



\textbf{Solution 3}

Note that the choice of a food item on a given day is symmetric, i.e. the number of ways to create the meal plan with a cake on Friday is the same as the number of ways to create the meal plan with pudding on Friday, and the same reasoning holds for the other desserts. Since every meal plan is counted by the summation of the $4$ aforementioned plans (note that Friday's dessert has to be one of the $4$ given desserts) and that these cases are mutually exclusive (i.e. you cannot make both a cake and pudding on Friday), each case results in a quarter of the total meal plans (Since there are $4$ desserts, we multiply by $\frac{1}{4}$). The total number of plans with no restrictions can be counted with constructive counting, as follows:

 We note that there are $4$ ways to choose the first dessert. Then, each dessert thereafter must be distinct from the prior one. Since there are $4$ options, and $1$ of them has been taken by the prior, each following dessert can be chosen in $(4 - 1) = 3$ ways. Thus, since there are $6$ desserts other than the first, the total number (without restrictions) is $4 \cdot 3^6$

 Thus, by our symmetry argument derived prior, we know that the number of desired meal plans is $\frac{\text{\# plans w/o restrictions}}{4} = \frac{4 \cdot 3^6}{4} = 3^6 = 729$, choice $\boxed{A}$




\textbf{Solution 4 (Answer choices)}

Notice that if we consider Friday cake day, Saturday has $3$ choices. Thus, we consider everything from Sunday to Thursday. Suppose Sunday has $4$ choices, and all the rest of the days have $3$ choices. If we multiply $4 \times 3^5$, we obtain $972$. However, if Wednesday had not been a cake day, there would only be $2$ choices for Thursday. Hence, we know that the answer must be less than $972$. Thus, $729 \Rightarrow \boxed{A}$.

 


	\item (\textit{AMC 12 12B 2012}) It takes Clea 60 seconds to walk down an escalator when it is not moving, and 24 seconds when it is moving. How many seconds would it take Clea to ride the escalator down when she is not walking?


 $\textbf{(A)}\ 36\qquad\textbf{(B)}\ 40\qquad\textbf{(C)}\ 42\qquad\textbf{(D)}\ 48\qquad\textbf{(E)}\ 52$

 \textbf{Answer: }B



\textbf{Solution 1}

She walks at a rate of $x$ units per second to travel a distance $y$. As $vt=d$, we find $60x=y$ and $24*(x+k)=y$, where $k$ is the speed of the escalator. Setting the two equations equal to each other, $60x=24x+24k$, which means that $k=1.5x$. Now we divide $60$ by $1.5$ because you add the speed of the escalator but remove the walking, leaving the final answer that it takes to ride the escalator alone as $\boxed{\textbf{(B)}\ 40}$




\textbf{Solution 2}

We write two equations using distance = rate * time. Let $r$ be the rate she is walking, and $e$ be the speed the escalator moves. WLOG, let the distance of the escalator be $120$, as the distance is constant. Thus, our $2$ equations are $120 = 60r$ and $120 = 24(r+e)$. Solving for $e$, we get $e = 3$. Thus, it will take Clea $\dfrac{120}{3} = \boxed{\textbf{(B)}\ 40}$ seconds.

 



\textbf{Solution 3}

Clea covers $\dfrac{1}{60}$ of the escalator every second. Say the escalator covers $\dfrac{1}{r}$ of the escalator every second. Since Clea and the escalator cover the entire escalator in $24$ seconds, we can use distance $=$ rate $\cdot$ time to get $24\left(\dfrac{1}{60} + \dfrac{1}{r}\right) = 1$. Solving gives us $r = 40$, so if Clea were to just stand on the escalator, it would take her $\boxed{\textbf{(B)}\ 40}$ seconds to get down. ~Extremelysupercooldude

 


	\item (\textit{AMC 12 12B 2012}) What is the area of the polygon whose vertices are the points of intersection of the curves $x^2 + y^2 =25$ and $(x-4)^2 + 9y^2 = 81$?


 $\textbf{(A)}\ 24\qquad\textbf{(B)}\ 27\qquad\textbf{(C)}\ 36\qquad\textbf{(D)}\ 37.5\qquad\textbf{(E)}\ 42$

 \textbf{Answer: }B



\textbf{Solution 1}

The first curve is a circle with radius $5$ centered at the origin, and the second curve is an ellipse with center $(4,0)$ and end points of $(-5,0)$ and $(13,0)$. Finding points of intersection, we get $(-5,0)$, $(4,3)$, and $(4,-3)$, forming a triangle with height of $9$ and base of $6.$ So the area of this triangle is $9 \cdot 6 \cdot 0.5 =\boxed{27 \textbf{ (B)}}.$




\textbf{Solution 2}

Given the equations $x^2 + y^2 =25$ and $(x-4)^2 + 9y^2 = 81$, we can substitute $y^2=25-x^2$ from the first equation and plug it in to the 2nd equation, giving us $(x-4)^2+9(25-x^2)=81$. After rearranging, $8x^2+8x-160=0$ or $x^2+x-20=0$. The solutions are $x=-5$ and $x=4$. This gives us the points $(-5,0),(4,3)$,and $(4,-3)$. The area of the triangle formed by these points is $27=\fbox{B}$




\textbf{Solution 3}

Using the algebra from solution 2, we find the points of intersection to be $(-5,0),(4,3)$,and $(4,-3)$. Using the shoelace theorem, we can easily find the area of the triangle to be $27=\fbox{B}$.

 More information on the shoelace theorem:\href{https://artofproblemsolving.com/wiki/index.php/Shoelace\_Theorem}{https://artofproblemsolving.com/wiki/index.php/Shoelace\_Theorem}

 


	\item (\textit{AMC 12 12B 2012}) In the equation below, $A$ and $B$ are consecutive positive integers, and $A$, $B$, and $A+B$ represent number bases: 
 \[132_A+43_B=69_{A+B}.\]
What is $A+B$?


 $\textbf{(A)}\ 9\qquad\textbf{(B)}\ 11\qquad\textbf{(C)}\ 13\qquad\textbf{(D)}\ 15\qquad\textbf{(E)}\ 17$

 \textbf{Answer: }C



\textbf{Solution 1}

Change the equation to base 10: 
 \[A^2 + 3A +2 + 4B +3= 6A + 6B + 9\] 
 \[A^2 - 3A - 2B - 4=0\] 

 Either $B = A + 1$ or $B = A - 1$, so either $A^2 - 5A - 6, B = A + 1$ or $A^2 - 5A - 2, B = A - 1$. The second case has no integer roots, and the first can be re-expressed as $(A-6)(A+1) = 0, B = A + 1$. Since A must be positive, $A = 6, B = 7$ and $A+B = \boxed{\textbf{(C)}\ 13}$




\textbf{Solution 2 (Answer Choices)}

We can eliminate answer choice $\textbf{(A)}$ because you can't have a $9$ in base $9$. Now we know that A and B are consecutive, so we can just test answers. You will only have to test at most $8$ cases. Eventually, after testing a few cases, you will find that $A=6$ and $B=7$. The solution is $\boxed{\mathbf{(C)}\ 13}$.

 


	\item (\textit{AMC 12 12B 2012}) How many sequences of zeros and ones of length 20 have all the zeros consecutive, or all the ones consecutive, or both?


 $\textbf{(A)}\ 190\qquad\textbf{(B)}\ 192\qquad\textbf{(C)}\ 211\qquad\textbf{(D)}\ 380\qquad\textbf{(E)}\ 382$

 \textbf{Answer: }D or E (both were accepted)



\textbf{Solution 1}

There are $\binom{20}{2}$ selections; however, we count these twice, therefore

 $2\cdot\binom{20}{2} = \boxed{\textbf{(D)}\ 380}$. The wording of the question implies D, not E.

 However, MAA decided to accept both D and E.

 



\textbf{Solution 2}

Consider the 20 term sequence of $0$'s and $1$'s.  Keeping all other terms 1, a sequence of $k>0$ consecutive 0's can be placed in $21-k$ locations. That is, there are 20 strings with 1 zero, 19 strings with 2 consecutive zeros, 18 strings with 3 consecutive zeros, ..., 1 string with 20 consecutive zeros. Hence there are $20+19+\cdots+1=\binom{21}{2}$ strings with consecutive zeros. The same argument shows there are $\binom{21}{2}$ strings with consecutive 1's. This yields $2\binom{21}{2}$ strings in all. However, we have counted twice those strings in which all the 1's and all the 0's are consecutive. These are the cases $01111...$, $00111...$, $000111...$, ..., $000...0001$ (of which there are 19) as well as the cases $10000...$, $11000...$, $111000...$, ..., $111...110$ (of which there are 19 as well). This yields $2\binom{21}{2}-2\cdot19=\boxed{\textbf{(E)}\ 382}$




\textbf{Solution 3}

First, we think of ways to make all the $1$'s consecutive. If there are no consecutive $1$'s, there are $\binom{20}{0}$ ways to order them. If there is one consecutive $1$, there are $\binom{20}{1}$ ways to order them. If there are two consecutive $1$'s, then there are $\binom{19}{1}$ ways to order them (We treat the two $1$'s like a block, and then order that block with 18 other $0$'s). Continuing in this fashion, there are $\binom{20}{0} + \binom{20}{1} + \binom{19}{1} + \cdots + \binom{1}{1} = 1 + 20 + 19 + \cdots + 2 + 1 = 210 + 1 = 211$ ways to order consecutive $1$'s. From symmetry, there are also $211$ ways to order the $0$'s. Now, from PIE, we subtract out the cases where both the $1$'s and the $0$'s are consecutive. We do this because when counting the ways to order the $1$'s, we counted all of these cases once. Then, we did so again when ordering the $0$'s. So, to only have all of these cases once, we must subtract them. If $1$ is the leftmost digit, then there are $20$ cases where all the $1$'s and $0$'s are consecutive (we basically are choosing how many $1$'s are consecutive, and there are $20$ possibilities. All other digits become $0$, which are automatically consecutive since the $1$'s are consecutive. There are also $20$ cases when $0$ is the left-most digit. Thus, there are a total of $211 + 211 - 20 - 20 = \boxed{\textbf{(E)}\ 382}$. But, from the way the problem is worded, it somewhat implies that the orderings must include both $1$'s and $0$'s, so the answer would then be $\boxed{\textbf{(D)}\ 380}$ after we subtract out the cases where the orderings are either all $1$'s or all $0$'s. But, since this is unclear, MAA accepted both $\boxed{\textbf{(D}\ 380}$ and $\boxed{\textbf{(E)}\ 382}$ as acceptable answers.

 


 



\textbf{Solution 4}

We consider two cases, and subtract their overcount.

 
Case $1$: Consecutive $0$s 

 If we have one consecutive $0$, then we have $20$ ways.If we have two consecutive $0$s, then we have $19$ ways by thinking of the two consecutives as a block.Continuing this pattern, if we have twenty consecutive $0$s, then we have only $1$ way.

 Therefore, we have $20+19+\cdots+1=\binom{21}{2}$ ways for this case. 

 
Case $2$: Consecutive $1$s

 Notice that if we just swap every $0$ to a $1$ in the previous case, we also have a valid arrangement.Hence, we also have $20+19+\cdots+1=\binom{21}{2}$ ways for this case.

 
Overcount: Notice that we can have BOTH the $0$s and the $1$s be consecutive. These are the cases $01111...$, $00111...$, $000111...$, ..., $000...0000$ which gives us $20$ ways being overcounted. If we invert the $0$s to $1$s, we similarly have $20$ more ways, hence we need to subtract $40$ from our total count. (Note: this method of overcounting subtracts out the all $0$s and the all $1$s case since the problem implies that there needs to be at least one of each)

 
So we have $210 + 210 - 40 = 380$ ways which gives us $\boxed{\textbf{(D)}\ 380}$. 

 


 


	\item (\textit{AMC 12 12B 2012}) Two parabolas have equations $y= x^2 + ax +b$ and $y= x^2 + cx +d$, where $a$, $b$, $c$, and $d$ are integers, each chosen independently by rolling a fair six-sided die. What is the probability that the parabolas will have at least one point in common?


 $\textbf{(A)}\ \frac{1}{2}\qquad\textbf{(B)}\ \frac{25}{36}\qquad\textbf{(C)}\ \frac{5}{6}\qquad\textbf{(D)}\ \frac{31}{36}\qquad\textbf{(E)}\ 1$

 \textbf{Answer: }D



\textbf{Solution 1}

Set the two equations equal to each other: $x^2 + ax + b = x^2 + cx + d$. Now remove the x squared and get $x$'s on one side: $ax-cx=d-b$. Now factor $x$: $x(a-c)=d-b$. If $a$ cannot equal $c$, then there is always a solution, but if $a=c$, a $1$ in $6$ chance, leaving a $1080$ out $1296$, always having at least one point in common. And if $a=c$, then the only way for that to work, is if $d=b$, a $1$ in $36$ chance, however, this can occur $6$ ways, so a $1$ in $6$ chance of this happening. So adding one  thirty sixth to $\frac{1080}{1296}$, we get the simplified fraction of $\frac{31}{36}$; answer $\boxed{(D)}$.

 



\textbf{Solution 2}

Proceed as above to obtain $x(a-c)=d-b$. The probability that the parabolas have at least 1 point in common is 1 minus the probability that they do not intersect. The equation $x(a-c)=d-b$ has no solution if and only if $a=c$ and $d\neq b$. The probability that $a=c$ is $\frac{1}{6}$ while the probability that $d\neq b$ is $\frac{5}{6}$. Thus we have $1-\left(\frac{1}{6}\right)\left(\frac{5}{6}\right)=\frac{31}{36}$ for the probability that the parabolas intersect.

 



\textbf{Solution 3}

Clearly, $ax + b = cx + d$. Imagine the two sides as lines - they will have no solutions when the two lines are parallel (eg. have the same gradient) which is when $a$ is not equal to $c$. Also, if $b = d$ and $a = c$, they're the same line so we must add one case. There are $36$ combinations of $a$ and $c$, of which they are equal in $6$ - but we must subtract 1 as if $a=c$ but $b=d$ they still intersect and have solutions. So we subtract this to obtain $\frac{36}{36} - \frac{5}{36} = \frac{31}{36}$. 

 



\textbf{Solution 4}

$x^2+ax+b=(x+ \frac{a}{2})^2+b- \frac{a^2}{4}$
$x^2+cx+d=(x+ \frac{c}{2})^2+d- \frac{c^2}{4}$

 The only case where these two functions have no intersections is when the x-values of the turning point are the same but the y-values are not the same.

 $\therefore \frac{a}{2}= \frac{c}{2} \implies a=c \quad\quad probability = \frac{1}{6} \times \frac{1}{6} \times 6= \frac{1}{6}$
$b- \frac{a^2}{4} \neq d- \frac{c^2}{4} \implies b \neq d \quad\quad probability = 1- \frac{1}{6}\times \frac{1}{6}\times 6= \frac{5}{6}$
$\therefore 1- \frac{1}{6}\times \frac{5}{6} = 1- \frac{5}{36} = \frac{31}{36}$



	\item (\textit{AMC 12 12B 2012}) Bernardo and Silvia play the following game. An integer between $0$ and $999$ inclusive is selected and given to Bernardo. Whenever Bernardo receives a number, he doubles it and passes the result to Silvia. Whenever Silvia receives a number, she adds $50$ to it and passes the result to Bernardo. The winner is the last person who produces a number less than $1000$. Let $N$ be the smallest initial number that results in a win for Bernardo. What is the sum of the digits of $N$?


 $\textbf{(A)}\ 7\qquad\textbf{(B)}\ 8\qquad\textbf{(C)}\ 9\qquad\textbf{(D)}\ 10\qquad\textbf{(E)}\ 11$

 \textbf{Answer: }A



\textbf{Solution 1}

The last number that Bernardo says has to be between 950 and 999. Note that $1\rightarrow 2\rightarrow 52\rightarrow 104\rightarrow 154\rightarrow 308\rightarrow 358\rightarrow 716\rightarrow 766$ contains 4 doubling actions. Thus, we have $x \rightarrow 2x \rightarrow 2x+50 \rightarrow 4x+100 \rightarrow 4x+150 \rightarrow 8x+300 \rightarrow 8x+350 \rightarrow 16x+700$. 

 Thus, $950<16x+700<1000$. Then, $16x>250 \implies x \geq 16$. 

 Because we are looking for the smallest integer $x$, $x=16$. Our answer is $1+6=\boxed{7}$, which is A.

 



\textbf{Solution 2}

Work backwards. The last number Bernardo produces must be in the range $[950,999]$. That means that before this, Silvia must produce a number in the range $[475,499]$. Before this, Bernardo must produce a number in the range $[425,449]$. Before this, Silvia must produce a number in the range $[213,224]$. Before this, Bernardo must produce a number in the range $[163,174]$. Before this, Silvia must produce a number in the range $[82,87]$. Before this, Bernardo must produce a number in the range $[32,37]$. Before this, Silvia must produce a number in the range $[16,18]$. Silvia could not have added 50 to any number before this to obtain a number in the range $[16,18]$, hence the minimum $N$ is 16 with the sum of digits being $\boxed{\textbf{(A)}\ 7}$.

 



\textbf{Solution 3}

If our first number is $N,$ then the sequence of numbers will be 
 \[2N,~2N+50,~4N+100,~4N+150,~8N+300,~8N+350,~16N+700,~16N+750\] Note that we cannot go any further because doubling gives an extra $1500$ at the end, which is already greater than $1000.$ The smallest $N$ will be given if $16N+750>1000>16N+700 \implies 15<N<19.$ Since the problem asks for the smallest possible value of $N,$ we get $16,$ and the sum of its digits is $1+6=\boxed{\textbf{(A)}~7}$




\textbf{Solution 4}

If Bernardo is to win, he must produce a number $n$ such that $n \geq 950$. This would mean Silvia would have to hand him a number $r$ such that $r \geq \frac{950}{2}=475$. This means Bernardo would have to give Silvia a number $s$ such that $s \geq 475-50=425$. Repeating this process gives the following sequence: 
 \[B=213, S=163, B=82, S=32, B=16\] Thus, the smallest possible value is $16$ and the answer is $1+6=\boxed{\textbf{(A)}7}$

            \begin{itemize}
        \item Note: For all the values of $B$ and $S$, we always round up a value if the inequality gives a fraction, since we're using the $\geq$                inequality.
            \end{itemize}
            


	\item (\textit{AMC 12 12B 2012}) Jesse cuts a circular paper disk of radius $12$ along two radii to form two sectors, the smaller having a central angle of $120$ degrees. He makes two circular cones, using each sector to form the lateral surface of a cone. What is the ratio of the volume of the smaller cone to that of the larger one?


 $\textbf{(A)}\ \frac{1}{8}\qquad\textbf{(B)}\ \frac{1}{4}\qquad\textbf{(C)}\ \frac{\sqrt{10}}{10}\qquad\textbf{(D)}\ \frac{\sqrt{5}}{6}\qquad\textbf{(E)}\ \frac{\sqrt{5}}{5}$

 \textbf{Answer: }C



\textbf{Solution 1}

Let us take the example of the smaller sector: notice that when it is folded into a cone, the radius becomes the slant height, and the circumference of the base of the cone is equal to the length of the sector. This gives us the slant height and radius of the cone's base (which comes from the circumference), and if we want to get the height of the cone, we simply have to do $s^2 - r_b^2$ by the Pythagorean theorem, where $s$ is the slant height and $r_b$ is the base of the radius. Repeat this process for the larger cone. It is easy to derive that the ratio of the volume of the small cone to that of the large one is simply $\frac{r_s^2h_s}{r_l^2h_l}$. 

 This gives us the answer of $\boxed{\mathbf{C}}$. 

 



\textbf{Solution 2}

If the original radius is $12$, then the circumference is $24\pi$; since arcs are defined by the central angles, the smaller arc, $120$ degree angle, is half the size of the larger sector. so the smaller arc is $8\pi$, and the larger is $16\pi$. Those two arc lengths become the two circumferences of the new cones; so the radius of the smaller cone is $4$ and the larger cone is $8$. Using the Pythagorean theorem, the height of the larger cone is $4\cdot\sqrt{5}$ and the smaller cone is $8\cdot\sqrt{2}$, and now for volume just square the radii and multiply by $\frac{1}{3}$ of the height to get the volume of each cone: $128\cdot\sqrt{2}$ and $256\cdot\sqrt{5}$ [both multiplied by three as ratio come out the same. now divide the volumes by each other to get the final ratio of $\boxed{\textbf{(C) } \frac{\sqrt{10}}{10}}$



	\item (\textit{AMC 12 12B 2012}) Amy, Beth, and Jo listen to four different songs and discuss which ones they like. No song is liked by all three. Furthermore, for each of the three pairs of the girls, there is at least one song liked by those girls but disliked by the third. In how many different ways is this possible?


 $\textbf{(A)}\ 108\qquad\textbf{(B)}\ 132\qquad\textbf{(C)}\ 671\qquad\textbf{(D)}\ 846\qquad\textbf{(E)}\ 1105$

 \textbf{Answer: }B



\textbf{Solution 1}

Let the ordered triple $(a,b,c)$ denote that $a$ songs are liked by Amy and Beth, $b$ songs by Beth and Jo, and $c$ songs by Jo and Amy. The only possible triples are $(1,1,1), (2,1,1), (1,2,1)(1,1,2)$. 

 To show this, first observe these are all valid conditions. Second, note that none of $a,b,c$ can be bigger than 3. Suppose otherwise, that $a = 3$. Without loss of generality, say that Amy and Beth like songs 1, 2, and 3. Then because there is at least one song liked by each pair of girls, we require either $b$ or $c$ to be at least 1. In fact, we require either $b$ or $c$ to equal 1, otherwise there will be a song liked by all three. Suppose $b = 1$. Then we must have $c=0$ since no song is liked by all three girls, which contradicts with the statement "for each of the three pairs of the girls, there is at least one song liked by those two girls".

 \textbf{Case 1}: How many ways are there for $(a,b,c)$ to equal $(1,1,1)$? There are 4 choices for which song is liked by Amy and Beth, 3 choices for which song is liked by Beth and Jo, and 2 choices for which song is liked by Jo and Amy. The fourth song can be liked by only one of the girls, or none of the girls, for a total of 4 choices. So $(a,b,c)=(1,1,1)$ in $4\cdot3\cdot2\cdot4 = 96$ ways.

 \textbf{Case 2}: To find the number of ways for $(a,b,c) = (2,1,1)$ (without order), observe there are $\binom{4}{2} = 6$ choices of songs for the first pair of girls. There remain 2 choices of songs for the next pair (who only like one song). The last song is given to the last pair of girls. But observe that we can let any three pairs of the girls to like two songs, so we multiply by 3. In this case there are $(6\cdot2)\cdot3=36$ ways for the girls to like the songs.

 That gives a total of $96 + 36 = 132$ ways for the girls to like the songs, so the answer is $\boxed{(\textrm{\textbf{B}})}$.

 



\textbf{Solution 2}

Let $AB, BJ$, and $AJ$ denote a song that is liked by Amy and Beth (but not Jo), Beth and Jo (but not Amy), and Amy and Jo (but not Beth), respectively. Similarly, let $A, B, J,$ and $N$ denote a song that is liked by only Amy, only Beth, only Jo, and none of them, respectively. Since we know that there is at least $1\: AB, BJ$, and $AJ$, they must be $3$ songs out of the $4$ that Amy, Beth, and Jo listened to. The fourth song can be of any type $N, A, B, J, AB, BJ$, and $AJ$ (there is no $ABJ$ because no song is liked by all three, as stated in the problem.) Therefore, we must find the number of ways to rearrange $AB, BJ, AJ$, and a song from the set $\{N, A, B, J, AB, BJ, AJ\}$.

 Case 1: Fourth song = $N, A, B, J$

 Note that in Case 1, all four of the choices for the fourth song are different from the first three songs.

 Number of ways to rearrange = $(4!)$ rearrangements for each choice $*\: 4$ choices = $96$.

 Case 2: Fourth song = $AB, BJ, AJ$

 Note that in Case $2$, all three of the choices for the fourth song repeat somewhere in the first three songs.

 Number of ways to rearrange = $(4!/2!)$ rearrangements for each choice $*\: 3$ choices = $36$.

 $96 + 36 = \boxed{\textbf{(B)} \: 132}$.

 



\textbf{Solution 3}

There are $\binom{4}{3}$ ways to choose the three songs that are liked by the three pairs of girls.

 There are $3!$ ways to determine how the three songs are liked, or which song is liked by which pair of girls.

 In total, there are $\binom{4}{3}\cdot3!$ possibilities for the first $3$ songs.

 There are $3$ cases for the 4th song, call it song D.

 Case $1$: D is disliked by all $3$ girls $\implies$ there is only $1$ possibility.

 Case $2$: D is liked by exactly $1$ girl $\implies$ there are $3$ possibility.

 Case $3$: D is liked by exactly $2$ girls $\implies$ there are $3$ pairs of girls to choose from. However, there's overlap when the other song liked by the same pair of girl is counted as the 4th song at some point, in which case D would be counted as one of the first $3$ songs liked by the same girls.

 Counting the overlaps, there are $3$ ways to choose the pair with overlaps and $4\cdot3=12$ ways to choose what the other $2$ pairs like independently. In total, there are $3\cdot12=36$ overlapped possibilities.

 Finally, there are $\binom{4}{3}\cdot3!\cdot(3+1+3)-36=132$ ways for the songs to be likely by the girls. $\boxed{\mathrm{(B)}}$




\textbf{Solution 4}

This is a bipartite graph problem, with the girls as left vertices and songs as right vertices. An edge connecting left vertex and right vertex means that a girl likes a song.

 Condition 1: "No song is liked by all three", means that the degree of right vertices is at most 2. 

 Condition 2: "for each of the three pairs of the girls, there is at least one song liked by those two girls but disliked by the third", means that for any pair of left vertices, there is at least a right vertex connecting to them.

 To meet condition 2, there are at least 3 right vertices with 2 edges connecting to left vertices.  There are 2 cases:

 Case 1: there are only 3 such right vertices. There are $\binom{4}{3}$ such vertices, with $3!$ ways of connections to the left vertices, total arrangements are $\binom{4}{3}\cdot3! = 24$. The fourth right vertex either has no edge to the 3 left vertices, or 1 edge to 1 of the 3 left vertices.  So there are $24\cdot(1+3) = 96$ ways.

 Case 2: there are 4 such right vertices, 2 of them have edges to the same pair of left vertices.  There are $\binom{4}{2}$ such vertices, with $3!$ ways of connections.  So there are $\binom{4}{2}\cdot3! = 36$ ways.

 Total ways are $96+36=132$.

 Another way is to overcount then subtract overlap ways.  Similar to previous case 1,  the fourth right vertex could have all possible connection to the left vertices except connecting to all 3, so it is $2^3-1=7$ ways, so the total ways are $\binom{4}{3}\cdot3!\cdot7 = 24\cdot7 = 168$.  But this overcounts the case 2 with 36 ways.  So total ways are $168-36=132$.

 -\href{https://artofproblemsolving.com/wiki/index.php/User:Junche}{junche}

 



\textbf{Solution 5}

Let's assign to each of the 3 pairs a song that they like, and that song is disliked by the other girl(satisfy the problem's 2nd condition). The first pair has 4 choice, 2nd has 3 choice, and last one has 2 choice. $4 \cdot 3 \cdot 2$. Then the last song can be liked/disliked freely, so there is $2^3 = 8$ total ways they can like/dislike the songs. Of these 8 cases, one of them is when they all like the song (which isn't allowed by the problem), so we subtract it away to have 7 cases. $4 \cdot 3 \cdot 2 \cdot 7 = 168$. But, we have overcounts when one pair of girls likes 2 songs, and the other girl dislikes those 2 songs. We can find the amount of cases that this happens by giving 1 pair of girls 2 songs that they like (and the other dislikes), and the other 2 pairs of girls 1 song that they like. There are 3 pairs of girls that we can choose to like 2 songs, so there are $3 \cdot \left(\binom {4}{2} \cdot 2 \cdot 1\right) = 36$ ways. $168 - 36 = \boxed {\textbf{(B) 132}}$.

 


	\item (\textit{AMC 12 12B 2012}) Square $PQRS$ lies in the first quadrant. Points $(3,0), (5,0), (7,0),$ and $(13,0)$ lie on lines $SP, RQ, PQ,$ and $SR$, respectively. What is the sum of the coordinates of the center of the square $PQRS$?


 $\textbf{(A)}\ 6\qquad\textbf{(B)}\ \frac{31}{5}\qquad\textbf{(C)}\ \frac{32}{5}\qquad\textbf{(D)}\ \frac{33}{5}\qquad\textbf{(E)}\ \frac{34}{5}$

 \textbf{Answer: }C



\textbf{Solution 1}


                    \begin{figure}[H]
                        \centering
                        \includegraphics[width=0.5\textwidth]{Solution/AMC_12_12B_2012_17_Solution_Figure_1.png}
                    \end{figure}
                    

 Construct the midpoints $E=(4,0)$ and $F=(10,0)$ and triangle $\triangle EMF$ as in the diagram, where $M$ is the center of square $PQRS$. Also construct points $G$ and $H$ as in the diagram so that $BG\parallel PQ$ and $CH\parallel QR$.

 Observe that $\triangle AGB\sim\triangle CHD$ while $PQRS$ being a square implies that $GB=CH$. Furthermore, $CD=6=3\cdot AB$, so $\triangle CHD$ is 3 times bigger than $\triangle AGB$. Therefore, $HD=3\cdot GB=3\cdot HC$. In other words, the longer leg is 3 times the shorter leg in any triangle similar to $\triangle AGB$.

 Let $K$ be the foot of the perpendicular from $M$ to $EF$, and let $x=EK$. Triangles $\triangle EKM$ and $\triangle MKF$, being similar to $\triangle AGB$, also have legs in a 1:3 ratio, therefore, $MK=3x$ and $KF=9x$, so $10x=EF=6$. It follows that $EK=0.6$ and $MK=1.8$, so the coordinates of $M$ are $(4+0.6,1.8)=(4.6,1.8)$ and so our answer is $4.6+1.8 = 6.4 =$ $\boxed{\mathbf{(C)}\ 32/5}$.

 


 



\textbf{Solution 2}


                    \begin{figure}[H]
                        \centering
                        \includegraphics[width=0.5\textwidth]{Solution/AMC_12_12B_2012_17_Solution_Figure_2.png}
                    \end{figure}
                    

 Let the four points be labeled $P_1$, $P_2$, $P_3$, and $P_4$, respectively. Let the lines that go through each point be labeled $L_1$, $L_2$, $L_3$, and $L_4$, respectively. Since $L_1$ and $L_2$ go through $SP$ and $RQ$, respectively, and $SP$ and $RQ$ are opposite sides of the square, we can say that $L_1$ and $L_2$ are parallel with slope $m$. Similarly, $L_3$ and $L_4$ have slope $-\frac{1}{m}$. Also, note that since square $PQRS$ lies in the first quadrant, $L_1$ and $L_2$ must have a positive slope. Using the point-slope form, we can now find the equations of all four lines: $L_1: y = m(x-3)$, $L_2: y = m(x-5)$, $L_3: y = -\frac{1}{m}(x-7)$, $L_4: y = -\frac{1}{m}(x-13)$.

 
Since $PQRS$ is a square, it follows that $\Delta x$ between points $P$ and $Q$ is equal to $\Delta y$ between points $Q$ and $R$. Our approach will be to find $\Delta x$ and $\Delta y$ in terms of $m$ and equate the two to solve for $m$. $L_1$ and $L_3$ intersect at point $P$. Setting the equations for $L_1$ and $L_3$ equal to each other and solving for $x$, we find that they intersect at $x = \frac{3m^2 + 7}{m^2 + 1}$. $L_2$ and $L_3$ intersect at point $Q$. Intersecting the two equations, the $x$-coordinate of point $Q$ is found to be $x = \frac{5m^2 + 7}{m^2 + 1}$. Subtracting the two, we get $\Delta x = \frac{2m^2}{m^2 + 1}$. Substituting the $x$-coordinate for point $Q$ found above into the equation for $L_2$, we find that the $y$-coordinate of point $Q$ is $y = \frac{2m}{m^2+1}$. $L_2$ and $L_4$ intersect at point $R$. Intersecting the two equations, the $y$-coordinate of point $R$ is found to be $y = \frac{8m}{m^2 + 1}$. Subtracting the two, we get $\Delta y = \frac{6m}{m^2 + 1}$. Equating $\Delta x$ and $\Delta y$, we get $2m^2 = 6m$ which gives us $m = 3$. Finally, note that the line which goes though the midpoint of $P_1$ and $P_2$ with slope $3$ and the line which goes through the midpoint of $P_3$ and $P_4$ with slope $-\frac{1}{3}$ must intersect at at the center of the square. The equation of the line going through $(4,0)$ is given by $y = 3(x-4)$ and the equation of the line going through $(10,0)$ is $y = -\frac{1}{3}(x-10)$. Equating the two, we find that they intersect at $(4.6, 1.8)$. Adding the $x$ and $y$-coordinates, we get $6.4 = 32/5$. Thus, answer choice $\boxed{\textbf{(C)}}$ is correct.

 



\textbf{Solution 3}

Note that the center of the square lies along a line that has an $x-$intercept of $\frac{3+5}{2}=4$, and also along another line with $x-$intercept $\frac{7+13}{2}=10$. Since these 2 lines are parallel to the sides of the square, they are perpendicular (since the sides of a square are). Let $m$ be the slope of the first line. Then $-\frac{1}{m}$ is the slope of the second line. We may use the point-slope form for the equation of a line to write $l_1:y=m(x-4)$ and $l_2:y=-\frac{1}{m}(x-10)$. We easily calculate the intersection of these lines using substitution or elimination to obtain $\left(\frac{4m^2+10}{m^2+1},\frac{6m}{m^2+1}\right)$ as the center or the square. Let $\theta$ denote the (acute) angle formed by $l_1$ and the $x-$axis. Note that $\tan\theta=m$. Let $s$ denote the side length of the square. Then $\sin\theta=s/2$. On the other hand the acute angle formed by $l_2$ and the $x-$axis is $90-\theta$ so that $\cos\theta=\sin(90-\theta)=s/6$. Then $m=\tan\theta=3$. Substituting into $\left(\frac{4m^2+10}{m^2+1},\frac{6m}{m^2+1}\right)$ we obtain $\left(\frac{23}{5},\frac{9}{5}\right)$ so that the sum of the coordinates is $\frac{32}{5}=6.4$. Hence the answer is $\framebox{C}$.

 



\textbf{Solution 4 (Fast)}

Suppose

 
 \[SP: y=m(x-3)\]
 \[RQ: y=m(x-5)\]
 \[PQ: -my=x-7\]
 \[SR: -my=x-13\]

 where $m >0$.

 Recall that the distance between two parallel lines $Ax+By+C=0$ and $Ax+By+C_1=0$ is $|C-C_1|/\sqrt{A^2+B^2}$, we have distance between $SP$ and $RQ$ equals to $2m/\sqrt{1+m^2}$, and the distance between $PQ$ and $SR$ equals to $6/\sqrt{1+m^2}$. Equating them, we get $m=3$.

 Then, the center of the square is just the intersection between the following two "mid" lines:

 
 \[L_1: y=3(x-4)\]
 \[L_2: -3y = x-10\]

 The solution is $(4.6,1.8)$, so we get the answer $4.6+1.8=6.4$.  $\framebox{C}$.

 



\textbf{Solution 5 (Trigonometry)}

Using the diagram shown in Solution 1, we can set angle $BCQ$ as $\theta$. We know that $AB=2$ and $BC=2$. Now using $AA$

 similarity, we know that $\triangle BCQ\sim\triangle ACP$ in a $1:2$ ratio. Now we can see that $CQ=-2$$\cos\theta$, therefore, 

 meaning that $PQ=-2$$\cos\theta$. $PQRS$ is a square, so $QR=-2$$\cos\theta$. We also know that $QCHR$ is also a square since its 

 angles are $90^\circ$ and all of its sides are equal. Because squares $PQRS$ and $QCHR$ have equal side lengths, they are 

 congruent leading to the conclusion that side $CH=-2$$\cos\theta$. Since $PQRS$ is a square, lines $PQ$ and $SR$ are parallel 

 meaning that angle $CDH$ and angle $BCQ$ are congruent. We can easily calculate that the length of $CD=6$ and furthermore that 

 $CH=6$$\sin\theta$. Setting $6\sin\theta=-2\cos\theta$, we get that $\tan\theta=-1/3$. This means $-1/3$ is the slope of line $PQ$

 and the lines parallel to it. This is good news because we are dealing with easy numbers. We can solve for the coordinates of 

 points $E$ and $F$ because they are the midpoints. This will make solving for the center of square $PQRS$ easier. $E=(4,0)$ and 

 $F=(10,0)$. We know the slopes of lines $MF$ and $ME$, which are $-1/3$ and $3$ respectively. Now we can get the two equations.

 
 \[\left\{\begin{array}{l}y=-1/3x+10/3\\y=3x-12\end{array}\right.\]

 By solving: 

 $-1/3x+10/3=3x-12,$

 we find that $x=4.6$. Then plugging $x$ back into one of the first equations, we can find $y$ and the final coordinate turns out to be $(4.6,1.8)$. Summing up the values of $x$ and $y$, you get $4.6+1.8=6.4=32/5$.   $\boxed{\mathbf{(C)}\ 32/5}$.

 



\textbf{Solution 6}


                    \begin{figure}[H]
                        \centering
                        \includegraphics[width=0.5\textwidth]{Solution/AMC_12_12B_2012_17_Solution_Figure_1.png}
                    \end{figure}
                    
$SP: y = mx - 3m$, $RQ: y = mx-5m$, $PQ: y = -\frac{1}{m}x + \frac{7}{m}$, $SR: y = -\frac{1}{m}x + \frac{13}{m}$

 Let $SP = RP = PQ = SR = a$, $\angle GAB = \angle HCD = \theta$, and the slope of $SP$ be $m$.

 When the slope of $SP$ is $m$, the slope of $SR$ is $-\frac{1}{m}$, $\tan \theta = m$, $\cot \theta = -\frac{1}{m}$
$\sin \theta = \frac{GB}{AB} = \frac{a}{2}$, $\cos \theta = \frac{HC}{CD} = \frac{a}{6}$

 As $\tan \theta = \frac{\sin \theta}{\cos \theta} = \frac{\frac{a}{2}}{\frac{a}{6}} = 3$, $m = 3$
$SP: y = 3x - 9$, $RQ: y = 3x-15$, $PQ: y = -\frac{1}{3}x + \frac{7}{3}$, $SR: y = -\frac{1}{3}x + \frac{13}{3}$
$3x - 9 = -\frac{1}{3}x + \frac{13}{3}$, $x = 4$, $y = 3$, $S = (4, 3)$
$3x-15 = -\frac{1}{3}x + \frac{7}{3}$, $x = 5.2$, $y = 0.6$, $Q = (5.2, 0.6)$
$M = (4.6, 1.8)$, $4.6 + 1.8 = \boxed{\mathbf{(C)}\ 32/5}$




\textbf{Solution 7 (Pure Trig and Similarity)}


                    \begin{figure}[H]
                        \centering
                        \includegraphics[width=0.5\textwidth]{Solution/AMC_12_12B_2012_17_Solution_Figure_4.png}
                    \end{figure}
                    Let $PQ=x$ and $\angle PCA=\theta.$ Draw the line $BE$ such that $E$ is on $AP$ and $BE\parallel PQ.$ Also, Draw the line $CF$ such that $F$ is on $DR$ and $CF\parallel RQ.$ Then $EB=FC=x$ and $\angle EBA=\angle FDC=\theta.$ Also, note that $AB=2$ and $CD=6.$ Hence:
 \[\cos(\theta)=\frac{x}{2},\sin(\theta)=\frac{x}{6}.\]Thus $\cos(\theta)=3\sin(\theta).$ Since $\theta<\frac{\pi}{2},$ $\sin(\theta)=\sqrt{1-\cos^2(\theta)},$ so $\cos(\theta)=3\sqrt{1-\cos^2(\theta)}.$ Hence, $\cos(\theta)=\frac{3}{\sqrt{10}}$ and $x=\frac{6}{\sqrt{10}}.$ Draw the perpendicular lines $EF_1\perp AD,PF_2\perp AD,RF_3\perp AD.$ Note that:
 \[F_1B=BE\cdot\cos(\theta)=\frac{6}{\sqrt{10}}\cdot\frac{3}{\sqrt{10}}=\frac{9}{5}.\]Hence:
 \[EF_1=\sqrt{EB^2-EF_1^2}=\sqrt{\frac{18}{5}-\frac{81}{25}}=\frac{3}{5}.\]Note that $\triangle EF_1B\sim\triangle PF_2C,$ so:
 \[\frac{PF_2}{EF_1}=\frac{F_2C}{F_1B}=\frac{AC}{AB}=2.\]Hence:
 \[PF_2=\frac{6}{5}, F_2C=\frac{18}{5}.\]So $P$ has coordinates:
 \[\left(7-\frac{18}{5},\frac{6}{5}\right)=\left(\frac{17}{5},\frac{6}{5}\right).\]Also note that $\triangle EF_1B\sim\triangle RF_3D,$ so:
 \[\frac{RF_3}{EF_1}=\frac{F_3D}{F_1B}=\frac{BD}{AB}=4.\]Hence:
 \[RF_3=\frac{12}{5}, F_3D=\frac{36}{5}.\]So $R$ has coordinates:
 \[\left(13-\frac{36}{5},\frac{12}{5}\right)=\left(\frac{29}{5},\frac{12}{5}\right).\]Hence, the center of square $PQRS,$ which is also the midpoint of $PR,$ has coordinates:
 \[\left(\frac{\frac{17}{5}+\frac{29}{5}}{2},\frac{\frac{6}{5}+\frac{12}{5}}{2}\right)=\left(\frac{23}{5},\frac{9}{5}\right).\]We thus see that the answer is:
 \[\frac{23}{5}+\frac{9}{5}=\boxed{\text{(C)}\frac{32}{5}}.\]



	\item (\textit{AMC 12 12B 2012}) Let $(a_1,a_2, \dots ,a_{10})$ be a list of the first $10$ positive integers such that for each $2 \le i \le 10$ either $a_i+1$ or $a_i-1$ or both appear somewhere before $a_i$ in the list. How many such lists are there?


 $\textbf{(A)}\ 120\qquad\textbf{(B)}\ 512\qquad\textbf{(C)}\ 1024\qquad\textbf{(D)}\ 181,440\qquad\textbf{(E)}\ 362,880$

 \textbf{Answer: }B



\textbf{Solution 1}

Let $1\leq k\leq 10$. Assume that $a_1=k$. If $k<10$, the first number appear after $k$ that is greater than $k$ must be $k+1$, otherwise if it is any number $x$ larger than $k+1$, there will be neither $x-1$ nor $x+1$ appearing before $x$. Similarly, one can conclude that if $k+1<10$, the first number appear after $k+1$ that is larger than $k+1$ must be $k+2$, and so forth.

 On the other hand, if $k>1$, the first number appear after $k$ that is less than $k$ must be $k-1$, and then $k-2$, and so forth.

 To count the number of possibilities when $a_1=k$ is given, we set up $9$ spots after $k$, and assign $k-1$ of them to the numbers less than $k$ and the rest to the numbers greater than $k$. The number of ways in doing so is $9$ choose $k-1$.

 Therefore, when summing up the cases from $k=1$ to $10$, we get

 
 \[\binom{9}{0} + \binom{9}{1} + \cdots + \binom{9}{9} = 2^9=512 ...... \framebox{B}\]




\textbf{Solution 2}

This problem is worded awkwardly. More simply, it asks: “How many ways can you order numbers 1-10 so that each number is one above or below some previous term?”

 
Then, the method becomes clear. For some initial number, WLOG examine the numbers greater than it. They always must appear in ascending order later in the list, though not necessarily as adjacent terms. Then, for some initial number, the number of possible lists is just the number of combination where this number of terms can be placed in 9 slots. For 9, that’s 1 number in 9 potential slots. For 8, that’s 2 numbers in 9 potential slots. 

 
 \[\binom{9}{0} + \binom{9}{1} + \cdots + \binom{9}{9} =512 \implies \boxed{B}\]




\textbf{Solution 3 (Noticing Stuff)}

If there is only 1 number, the number of ways to order would be 1.If there are 2 numbers, the number of ways to order would be 2.If there are 3 numbers, by listing out, the number of ways is 4.We can then make a conjecture that the problem is simply powers of 2. 

 \begin{verbatim}
$2^{10-1}=512$.
\end{verbatim}





\textbf{Solution 4 (fast and clever)}

Assume that the first element is any integer.Since we can only add the integer 1 more than the current max or 1 less the current min, there are 2 possibilities for each element after the first.Once we have created a set of 10 elements, we can shift all of the elements by some constant so that they will be the first 10 positive integers.Thus the total possibilities is $2^9$.

 



\textbf{Solution 5 (Recursion)}

For a list ${a_1, a_2, \dots, a_k}$ with $k$ terms, $2$ valid lists with $k+1$ terms can be created by $2$ ways:

 1. Add $a_{k+1}$ to the end of the list, making a new list ${a_1, a_2, \dots, a_k, a_{k+1} }$

 2. Increase the value of all existing terms by one, making a new list ${a_2, a_3, \dots, a_{k+1}}$. Then add $a_1$ to the end of the list, making a new list ${a_2, a_3, \dots, a_{k+1}, a_1}$

 Let $F(n)$ be the number of lists with $n$ elements, $F(n) = 2 \cdot F(n-1)$. As $F(2) = 2$, $F(n) = 2^{n-1}$, $F(10) = 2^9 = \boxed{\textbf{(B) } 512}$




\textbf{Solution 6 (Recursion but explained better)}

Let $(a_1,a_2 \dots a_k)$ be a valid list of $\{1,2 \dots k\}$. We wish to create a list of $\{1,2 \dots k+1\}$ using our valid list. 

 Claim: every $k-list$ must end with either $k$ or $1$. 

 Proof: Let last element is $2 \leq l \leq k-1$. Two cases:

 $l-1$ appears before $l$ in the list: We apply the same to $l-1$ and get $l-2$ appearing before in the list and so on. This means that the set $\{1,2 \dots l-1\}$ has exactly $k-1$ elements. This gives $(l-1) -1 + 1 = k-1 \implies l = k$. 

 $l+1$ appears before $l$ in the list: This tells us that $l+2$ appears before $l+1$, and $l+3$ before $l+2$ and so on. This means that the set $\{l+1, l+2 \dots l+j\}$ has exactly $k-1$ elements. This gives $l+j - (l+1) + 1 = k-1 \implies j = k-1 \implies l = 1$ (because $l+j$ is an element, and that is only possible if $l=1$). 

 Now, we construct our list of $\{1,2 \dots k+1 \}$. First case is to just put $k+1$ at the end of $(a_1,a_2 \dots a_k)$ which gives us every valid list ending with $k+1$. Now, we need all valid lists ending with $1$. Since our current list is of $\{1,2 \dots k\}$, if we add $1$ to every element, our list becomes a list of $\{2,3 \dots k+1 \}$. But note that our list is still a valid list. So, if we add $1$ to the end of this list, we will get every valid list of $k+1$ length ending with $1$. 

 Thus, from a valid list of $k$, you get 2 valid lists of $k+1$. Let $F(k)$ be lists of $k$. 
 \[F(k) = 2F(k-1), F(2) = 2 \implies F(k) = 2^{k-1} \implies \boxed{F(10) = 2^9 = 512}\]



	\item (\textit{AMC 12 12B 2012}) A unit cube has vertices $P_1,P_2,P_3,P_4,P_1',P_2',P_3',$ and $P_4'$. Vertices $P_2$, $P_3$, and $P_4$ are adjacent to $P_1$, and for $1\le i\le 4,$ vertices $P_i$ and $P_i'$ are opposite to each other. A regular octahedron has one vertex in each of the segments $P_1P_2$, $P_1P_3$, $P_1P_4$, $P_1'P_2'$, $P_1'P_3'$, and $P_1'P_4'$. What is the octahedron's side length?


 
                    \begin{figure}[H]
                        \centering
                        \includegraphics[width=0.5\textwidth]{Problem/AMC_12_12B_2012_19_Problem_Figure_1.png}
                    \end{figure}
                    
$\textbf{(A)}\ \frac{3\sqrt{2}}{4}\qquad\textbf{(B)}\ \frac{7\sqrt{6}}{16}\qquad\textbf{(C)}\ \frac{\sqrt{5}}{2}\qquad\textbf{(D)}\ \frac{2\sqrt{3}}{3}\qquad\textbf{(E)}\ \frac{\sqrt{6}}{2}$

 \textbf{Answer: }A



\textbf{Solution 1}


                    \begin{figure}[H]
                        \centering
                        \includegraphics[width=0.5\textwidth]{Solution/AMC_12_12B_2012_19_Solution_Figure_1.jpg}
                    \end{figure}
                    

 Observe the diagram above. Each dot represents one of the six vertices of the regular octahedron. Three dots have been placed exactly x units from the $(0,0,0)$ corner of the unit cube. The other three dots have been placed exactly x units from the $(1,1,1)$ corner of the unit cube. A red square has been drawn connecting four of the dots to provide perspective regarding the shape of the octahedron. Observe that the three dots that are near $(0,0,0)$ are each $(x)(\sqrt{2}$) from each other. The same is true for the three dots that are near $(1,1,1).$ There is a unique $x$ for which the rectangle drawn in red becomes a square. This will occur when the distance from $(x,0,0)$ to $(1,1-x, 1)$ is $(x)(\sqrt{2}$).

 
Using the distance formula we find the distance between the two points to be: $\sqrt{{(1-x)^2} + {(1-x)^2} + 1}$ = $\sqrt{2x^2 - 4x +3}$. Equating this to $(x)(\sqrt{2}$) and squaring both sides, we have the equation:

 $2{x^2} - 4x + 3$ = $2{x^2}$
$-4x + 3 = 0$
$x$ = $\frac{3} {4}$.

 
Since the length of each side is $(x)(\sqrt{2}$), we have a final result of $\frac{3 \sqrt{2}}{4}$. Thus, Answer choice $\boxed{\text{A}}$ is correct.

 
(If someone can draw a better diagram with the points labeled P1,P2, etc., I would appreciate it).

 --\href{https://artofproblemsolving.com/wiki/index.php?title=User:Jm314&action=edit&redlink=1}{Jm314} 14:55, 26 February 2012 (EST)

 



\textbf{Solution 2}

Standard 3D geometry, no coordinates. 

 
Let the tip of the octahedron on side $P_1P_3$ be $K_1$ and the opposite vertex be $K_2$. Our key is to examine the trapezoid $P_1K_1K_2P_3$. 

 
Let the side length of the octahedron be $s$. Then $P_1K_1 = \frac{s}{\sqrt{2}}$ and $P_3K_2 = 1 - \frac{2}{\sqrt{2}}$. Then, we have $P_1P_3 = \sqrt{2}$. Finally, we want to find $K_1K_2$, which is just double the height of half the octahedron. We can use Pythagorean Theorem to find that height as $\sqrt{2}s$. Now, we use the Pythagorean Theorem on the trapezoid. We get

 
 \[(\sqrt{2})^2 + (2\sqrt{2}-1)^2 = (s\sqrt{2})^2\] 
 \[s = \frac{3\sqrt{2}}{4}.\]




\textbf{Solution 3}


                    \begin{figure}[H]
                        \centering
                        \includegraphics[width=0.5\textwidth]{Solution/AMC_12_12B_2012_19_Solution_Figure_2.png}
                    \end{figure}
                    Let the length of $P_1A = a$, $P_1B = b$
$AB = a^2 + b^2$, $AB' = (1-b)^2 + (1-a)^2 + 1$, $AB = AB'$
$a^2 + b^2 = (1-b)^2 + (1-a)^2 + 1$, $a^2 + b^2 = 1 - 2b + b^2 + 1 - 2a + a^2 + 1$, $a+ b = \frac32$

 As $AC = BC$, $a^2 + P_1C^2 = b^2 + P_1C^2$, $a = b$, $a = \frac34$
$AB = \boxed{\textbf{(A) } \frac{3\sqrt{2}}{4}}$



	\item (\textit{AMC 12 12B 2012}) A trapezoid has side lengths $3$, $5$, $7$, and $11$. The sums of all the possible areas of the trapezoid can be written in the form of $r_1\sqrt{n_1}+r_2\sqrt{n_2}+r_3$, where $r_1$, $r_2$, and $r_3$ are rational numbers and $n_1$ and $n_2$ are positive integers not divisible by the square of any prime. What is the greatest integer less than or equal to $r_1+r_2+r_3+n_1+n_2$?


 $\textbf{(A)}\ 57\qquad\textbf{(B)}\ 59\qquad\textbf{(C)}\ 61\qquad\textbf{(D)}\ 63\qquad\textbf{(E)}\ 65$

 \textbf{Answer: }D



\textbf{Solution 1}

Name the trapezoid $ABCD$, where $AB$ is parallel to $CD$, $AB<CD$, and $AD<BC$. Draw a line through $B$ parallel to $AD$, crossing the side $CD$ at $E$. Then $BE=AD$, $EC=DC-DE=DC-AB$. One needs to guarantee that $BE+EC>BC$, so there are only three possible trapezoids:

 
 \[AB=3, BC=7, CD=11, DA=5, CE=8\]



	\item (\textit{AMC 12 12B 2012}) Square $AXYZ$ is inscribed in equiangular hexagon $ABCDEF$ with $X$ on $\overline{BC}$, $Y$ on $\overline{DE}$, and $Z$ on $\overline{EF}$. Suppose that $AB=40$, and $EF=41(\sqrt{3}-1)$. What is the side-length of the square?


 
                    \begin{figure}[H]
                        \centering
                        \includegraphics[width=0.5\textwidth]{Problem/AMC_12_12B_2012_21_Problem_Figure_1.png}
                    \end{figure}
                    
$\textbf{(A)}\ 29\sqrt{3} \qquad\textbf{(B)}\ \frac{21}{2}\sqrt{2}+\frac{41}{2}\sqrt{3}\qquad\textbf{(C)}\ 20\sqrt{3}+16$
$\textbf{(D)}\ 20\sqrt{2}+13\sqrt{3} \qquad\textbf{(E)}\ 21\sqrt{6}$

 \textbf{Answer: }A



\textbf{Solution 1}

We can, $\textsc{wlog}$, assume $Y$ coincides with $D$ and $CD\parallel AF$ as before. In which case, we will have $BC=EF=41(\sqrt{3}-1)$. So we have square $AXDZ$ inscribed in equiangular hexagon $ABCDEF$ with $X$ on $\overline{BC}$ and $Z$ on $\overline{EF}$.
                    \begin{figure}[H]
                        \centering
                        \includegraphics[width=0.5\textwidth]{Solution/AMC_12_12B_2012_21_Solution_Figure_1.png}
                    \end{figure}
                    Let $\angle BXA = \theta$; then $\angle BAX=60^\circ -\theta$. Let $BX=u$. In $\triangle ABX$ we have 
 \begin{align}     \frac{2s}{\sqrt{3}}=\frac{u}{\sin(60^\circ-\theta)}=\frac {40}{\sin\theta} \end{align}We also have $\angle CXD=90^\circ - \theta$ and $\angle CDX = \theta-30^\circ$. Let $CX=v$. In $\triangle CDX$ we have
 \begin{align}\tag{2}     \frac{2s}{\sqrt{3}}=\frac{v}{\sin(\theta-30^\circ)}=\frac {CD}{\cos\theta} \end{align}Now $BC=u+v=41(\sqrt{3}-1)$. From $(1)$ and $(2)$ we get
 \begin{align*}     41(\sqrt{3}-1) &= \frac{2s}{\sqrt{3}}\left(\sin(60^\circ-\theta)+\sin(\theta-30^\circ)\right) \\     &= \frac{2s}{\sqrt{3}} \cdot \frac{\sqrt{3}-1}2\cdot (\sin\theta + \cos\theta) \end{align*}
From $(1)$ we get $s\sin\theta = 20\sqrt{3}$ and therefore $s\cos\theta = \sqrt{s^2-3\cdot 20^2}$. Thus
 \[41(\sqrt{3}-1) = \frac{\sqrt{3}-1}{\sqrt{3}}(20\sqrt{3}+\sqrt{s^2-3\cdot 20^2})\]which simplifies to
 \[3\cdot 21^2  = s^2-3\cdot 20^2.\]Since $(20, 21, 29)$ is a Pythagorean triple, we get $s=29\sqrt{3}$, i.e. $\framebox{A}$.

 



\textbf{Solution 2}

Extend $AF$ and $YE$ so that they meet at $G$. Then $\angle FEG=\angle GFE=60^{\circ}$, so $\angle FGE=60^{\circ}$ and because $AB$ is parallel to $YE$. Also, since $AX$ is parallel and equal to $YZ$, we get $\angle BAX = \angle ZYE$, hence $\triangle ABX$ is congruent to $\triangle YEZ$. We now get $YE=AB=40$.

 Let $a_1=EY=40$, $a_2=AF$, and $a_3=EF$.

 Drop a perpendicular line from $A$ to the line of $EF$ that meets line $EF$ at $K$, and a perpendicular line from $Y$ to the line of $EF$ that meets $EF$ at $L$, then $\triangle AKZ$ is congruent to $\triangle ZLY$ since $\angle YZL$ is complementary to $\angle KZA$. Then we have the following equations:

 
 \[\frac{\sqrt{3}}{2}a_2 = AK=ZL = ZE+\frac{1}{2} a_1\]
 \[\frac{\sqrt{3}}{2}a_1 = YL =ZK = ZF+\frac{1}{2} a_2\]

 The sum of these two yields that 

 
 \[\frac{\sqrt{3}}{2}(a_1+a_2) = \frac{1}{2} (a_1+a_2) + ZE+ZF =  \frac{1}{2} (a_1+a_2) + EF\]
 \[\frac{\sqrt{3}-1}{2}(a_1+a_2) = 41(\sqrt{3}-1)\]
 \[a_1+a_2=82\]
 \[a_2=82-40=42.\]

 So, we can now use the law of cosines in $\triangle AGY$:

 
 \[2AZ^2 = AY^2 = AG^2 + YG^2 - 2AG\cdot YG \cdot \cos 60^{\circ}\]
 \[= (a_2+a_3)^2 + (a_1+a_3)^2 - (a_2+a_3)(a_1+a_3)\]
 \[= (41\sqrt{3}+1)^2 + (41\sqrt{3}-1)^2 - (41\sqrt{3}+1)(41\sqrt{3}-1)\]
 \[= 6 \cdot 41^2 + 2 - 3 \cdot 41^2 + 1 = 3 (41^2 + 1) = 3\cdot 1682\]
 \[AZ^2 = 3 \cdot 841 = 3 \cdot 29^2\]

 Therefore $AZ = 29\sqrt{3} ... \framebox{A}$




\textbf{Solution 3}

First, we want to angle chase. Set $<YXC$ equal to $a$ degrees. 

 Now the key idea is that you want to relate the numbers that you have. You know $\overline{AB} = 40$ and that $\overline{EZ} + \overline{ZF} = 41(\sqrt{3}-1)$. We proceed with the Law of Sines.

 Call the side length of the square x. Then we are going to set a constant k equal to $\frac{\sin 120^{\circ}}{x}$, and this is consistent for every triangle in the diagram because all the angles of the hexagon are equiangular (and so they are all $120^{\circ}$). 

 Then we get the following process:
 \[\frac{\sin(90-a)}{40} = k\]
 \[\cos a = 40k\]

 \[\frac{\sin(a-30)}{\overline{EZ}} = k\]
 \[\sin(a-30) = \overline{EZ}\cdot k\]
 \[\frac{\sin(60-a)}{\overline{ZF}} = k\]
 \[\sin(60-a) = \overline{ZF}\cdot k\]
 \[\sin(a-30) + \sin(60-a) = k\cdot 41(\sqrt{3}-1)\]

 And now expanding using our trig formulas, we get:
 \[(\sin a + \cos a)(\frac{\sqrt{3}-1}{2} = k\cdot 41(\sqrt{3}-1)\]
 \[\sin a + \cos a = 82k\]
 \[\sin a = 42k\]

 And so now we have a triangle where $\cos a = 40k$ and $\sin a = 42k$. Put them in a triangle where the hypotenuse is 1. Then, by the Pythagorean Theorem, we get:
 \[\sqrt{(40k)^2 + (42k)^2} = 1\]
 \[3364k^2 = 1\]
 \[k = \frac{1}{58}\]

 And since $k = \frac{\sin(120^{\circ})}{x}$, then:
 \[x = \frac{\sqrt{3}}{2}\cdot58\]
 \[x = \boxed{29\sqrt{3}}\]

 Solution by IronicNinja

 



\textbf{Solution 4}

Let $EZ = x$, $\angle XAB = \alpha$
$\angle BAX = \angle EYZ$, $AX = YZ$, $\angle ZEY = \angle XBA$, $\triangle BAX \cong \triangle EYZ$ by $AAS$, $BX = EZ = x$
$\angle AXB = 180^\circ - 120^\circ - \alpha = 60^\circ - \alpha$
$\frac{XB}{\sin \angle XAB} = \frac{AX}{\sin \angle ABX} = \frac{AB}{\angle AXB}$, $\frac{x}{\sin \alpha} = \frac{AX}{\sin 120^\circ} = \frac{40}{\sin (60^\circ - \alpha)}$
$\angle ZAF = 120^\circ - 90^\circ - \alpha = 30^\circ - \alpha$
$ZF = 41(\sqrt{3} - 1) - x$, $\frac{ZF}{\sin \angle ZAF} = \frac{AZ}{ \sin \angle ZFA}$, $\frac{41(\sqrt{3} - 1) - x}{\sin (30^\circ - \alpha)} = \frac{AZ}{ \sin 120^\circ}$
$\frac{x}{\sin \alpha} = \frac{41(\sqrt{3} - 1) - x}{\sin (30^\circ - \alpha)}= \frac{40}{\sin(60^\circ - \alpha)}$
$40 \cdot \sin \alpha = x(\sin 60^\circ \cos \alpha - \cos 60^\circ \sin \alpha)$ $(1)$
$x(\sin 30^\circ \cos \alpha - \cos 30^\circ \sin \alpha) = [41(\sqrt{3} - 1) - x] \sin \alpha$ $(2)$

 By simplifying $(1)$ we get, $\frac{\sqrt{3}}{2} \cdot x \cdot \cos \alpha - \frac{\sin \alpha}{2} \cdot x = 40 \cdot \sin \alpha$ $(3)$

 By simplifying $(2)$ we get, $\frac{\cos \alpha}{2} \cdot x + \frac{2 - \sqrt{3}}{2} \cdot \sin \alpha \cdot x = 41(\sqrt{3} - 1)\sin \alpha$ $(4)$

 By $\sqrt{3} \cdot (4) - (3)$ we get, $\frac{2\sqrt{3} - 3 +1}{2} \cdot \sin \alpha \cdot x = [41(3-\sqrt{3}) - 40] \sin \alpha$
$(\sqrt{3} - 1)x = 83 - 41\sqrt{3}$, $x = 21 \sqrt{3} - 20$

 By the law of cosine $AX = \sqrt{(21 \sqrt{3} - 20)^2 + 40^2 - 2 \cdot (21 \sqrt{3} - 20) 40 \cdot \cos 120^\circ} = \boxed{\textbf{(A) } 29 \sqrt{3}}$



	\item (\textit{AMC 12 12B 2012}) A bug travels from $A$ to $B$ along the segments in the hexagonal lattice pictured below. The segments marked with an arrow can be traveled only in the direction of the arrow, and the bug never travels the same segment more than once. How many different paths are there?


 
                    \begin{figure}[H]
                        \centering
                        \includegraphics[width=0.5\textwidth]{Problem/AMC_12_12B_2012_22_Problem_Figure_1.png}
                    \end{figure}
                    
$\textbf{(A)}\ 2112\qquad\textbf{(B)}\ 2304\qquad\textbf{(C)}\ 2368\qquad\textbf{(D)}\ 2384\qquad\textbf{(E)}\ 2400$

 \textbf{Answer: }E



\textbf{Solution 1}


                    \begin{figure}[H]
                        \centering
                        \includegraphics[width=0.5\textwidth]{Solution/AMC_12_12B_2012_22_Solution_Figure_1.png}
                    \end{figure}
                    

 We can count the number of ways to get from each set of only solid arrows to the next set; if we pass any set, we can't go to points behind that set. We break the diagram up like this to ensure that we don't overcount; otherwise, we would be able to return through the outlined white arrow to previous points, making the casework very difficult. Then, we find the total number of cases by multiplying the number of ways we can get from each of the sections to the next one.

 There is $1$ way to get to any of the red arrows. From the first (top) red arrow, there are $2$ ways to get to each of the first and the second (top 2) blue arrows; from the second (bottom) red arrow, there are $3$ ways to get to each of the first and the second blue arrows. So, by symmetry there are in total $5$ ways to get from any of the red arrows to any of the blue arrows. 

 From each of the first and second blue arrows, there are respectively $4$ ways to get to each of the first and the second green arrows; from each of the third and the fourth blue arrows, there are respectively $8$ ways to get to each of the first and the second green arrows. Therefore, as there are $5$ ways to go from the start to blue, and another $4+4+8+8$ ways to go from blue to green, there are $120$ ways to get to the green arrows.

 Finally, from each of the first and second green arrows, there are respectively $2$ ways to get to the first orange arrow; from each of the third and the fourth green arrows, there are $3$ ways to get to the first orange arrow. Therefore there are $120 \cdot (2+2+3+3) = 1200$ ways to get to each of the orange arrows, hence $2400$ ways to get to the point $B$.  $\boxed{\textbf{(E)}\ 2400}$




\textbf{Solution 2 (using the answer choices)}


                    \begin{figure}[H]
                        \centering
                        \includegraphics[width=0.5\textwidth]{Solution/AMC_12_12B_2012_22_Solution_Figure_2.png}
                    \end{figure}
                    

 For every blue arrow, there are $2\cdot 2=4$ ways to reach it without using the reverse arrow since the bug can choose any of $2$ red arrows to pass through and $2$ black arrows to pass through.  If the bug passes through the white arrow, the red arrow that the bug travels through must be the closest to the first black arrow. Otherwise, the bug will have to travel through both red segments, which is impossible because now there is no path to take after the bug emerges from the reverse arrow.  Similarly, with the blue segments, the second black arrow taken must be the one that is closest to the blue arrow that was taken.  Also, it is trivial that the two black arrows taken must be different.  Therefore, if the reverse arrow is taken, the blue arrow taken determines the entire path and there is $1$ path for every arrow.  Since the bug cannot return once it takes a blue arrow, the answer must be divisible by $5$.  $\boxed{\textbf{(E)}\ 2400}$ is the only answer that is.

 



\textbf{Solution 3: Casework}

The main thing to notice is that, if the bug does not go backwards, then every vertical set of arrows can be used, as the bug could travel straight downwards and then use any arrow to continue right.

 Note: The motivation is quite weird so follow my numbers as they start from the left (point A) and go right (point B).

 Case 1: Bug does not go backwards.

 The number of cases for this is just each vertical set of arrows multiplied to each other (if you don't understand where I'm coming from, try to understand where these numbers come from!)

 And so we have $2\cdot2\cdot4\cdot4\cdot4\cdot2\cdot2 = 2^{10}$ cases.

 Case 2: The bug goes backwards once, either at the first arrow or third arrow.

 Here, we have to count the fact that there is a horizontal midline that the bug could not cross, or otherwise it would be stepping on the same edge twice.Back on first arrow: $2\cdot1\cdot2\cdot4\cdot4\cdot2\cdot2 = 2^{8}$ cases.The third arrow is the same. 

 Case 3: The bug goes backwards once, at the second arrow.

 Same thing as above, except since there are 4 arrows in the vertical set (plus one for the backwards arrow), then the calculations are a bit different.

 We have $2\cdot2\cdot4\cdot1\cdot2\cdot4\cdot2\cdot2 = 2^{9}$ cases.

 Notice that the first and third back arrows decrease the number of cases by a factor of $2^2$ and the second back arrow decreases the number of cases by $2^1$. Hence,

 1st + 2nd = $2^7$

 2nd + 3rd = $2^7$

 1st + 3rd = $2^6$

 1st + 2nd + 3rd = $2^5$

 And so the number of cases in total is $2^{10} + 2^9 + 2^8 + 2^8 + 2^7 + 2^7 + 2^6 + 2^5 = 1024 + 512 + 256 + 256 + 128 + 128 + 64 + 32 = \boxed{2400}$



	\item (\textit{AMC 12 12B 2012}) Consider all polynomials of a complex variable, $P(z)=4z^4+az^3+bz^2+cz+d$, where $a,b,c,$ and $d$ are integers, $0\le d\le c\le b\le a\le 4$, and the polynomial has a zero $z_0$ with $|z_0|=1.$ What is the sum of all values $P(1)$ over all the polynomials with these properties?


 $\textbf{(A)}\ 84\qquad\textbf{(B)}\ 92\qquad\textbf{(C)}\ 100\qquad\textbf{(D)}\ 108\qquad\textbf{(E)}\ 120$

 \textbf{Answer: }B



\textbf{Solution 1}

Since $z_0$ is a root of $P$, and $P$ has integer coefficients, $z_0$ must be algebraic. Since $z_0$ is algebraic and lies on the unit circle, $z_0$ must be a root of unity (Comment: this is not true. See this link: \href{https://artofproblemsolving.comhttp://math.stackexchange.com/questions/4323/are-all-algebraic-integers-with-absolute-value-1-roots-of-unity}{[1]}). Since $P$ has degree 4, it seems reasonable (and we will assume this only temporarily) that $z_0$ must be a 2nd, 3rd, or 4th root of unity. These are among the set $\{\pm1,\pm i,(-1\pm i\sqrt{3})/2\}$. Since complex roots of polynomials come in conjugate pairs, we have that $P$ has one (or more) of the following factors: $z+1$, $z-1$, $z^2+1$, or $z^2+z+1$. If $z=1$ then $a+b+c+d+4=0$; a contradiction since $a,b,c,d$ are non-negative. On the other hand, suppose $z=-1$. Then $(a+c)-(b+d)=4$. This implies $a+b=8,7,6,5,4$ while $b+d=4,3,2,1,0$ correspondingly. After listing cases, the only such valid $a,b,c,d$ are $4,4,4,0$, $4,3,3,0$, $4,2,2,0$, $4,1,1,0$, and $4,0,0,0$. 

 Now suppose $z=i$. Then $4=(a-c)i+(b-d)$ whereupon $a=c$ and $b-d=4$. But then $a=b=c$ and $d=a-4$. This gives only the cases $a,b,c,d$ equals $4,4,4,0$, which we have already counted in a previous case.

 Suppose $z=-i$. Then $4=i(c-a)+(b-d)$ so that $a=c$ and $b=4+d$. This only gives rise to $a,b,c,d$ equal $4,4,4,0$ which we have previously counted. 

 Finally suppose $z^2+z+1$ divides $P$. Using polynomial division ((or that $z^3=1$ to make the same deductions) we ultimately obtain that $b=4+c$. This can only happen if $a,b,c,d$ is $4,4,0,0$. 

 Hence we've the polynomials
 \[4x^4+4x^3+4x^2+4x\]
 \[4x^4+4x^3+3x^2+3x\]
 \[4x^4+4x^3+2x^2+2x\]
 \[4x^4+4x^3+x^2+x\]
 \[4x^4+4x^3\]
 \[4x^4+4x^3+4x^2\]However, by inspection $4x^4+4x^3+4x^2+4x+4$ has roots on the unit circle, because $x^4+x^3+x^2+x+1=(x^5-1)/(x-1)$ which brings the sum to 92 (choice B).  Note that this polynomial has a 5th root of unity as a root. We will show that we were \textit{almost} correct in our initial assumption; that is that $z_0$ is at most a 5th root of unity, and that the last polynomial we obtained is the last polynomial with the given properties. Suppose that $z_0$ in an $n$th root of unity where $n>5$, and $z_0$ is not a 3rd or 4th root of unity. (Note that 1st and 2nd roots of unity are themselves 4th roots of unity). If $n$ is prime, then \textit{every} $n$th root of unity except 1 must satisfy our polynomial, but since $n>5$ and the degree of our polynomial is 4, this is impossible. Suppose $n$ is composite. If it has a prime factor $p$ greater than 5 then again every $p$th root of unity must satisfy our polynomial and we arrive at the same contradiction. Therefore suppose $n$ is divisible only by 2,3,or 5. Since by hypothesis $z_0$ is not a 2nd or 3rd root of unity, $z_0$ must be a 5th root of unity. Since 5 is prime, every 5th root of unity except 1 must satisfy our polynomial. That is, the other 4 complex 5th roots of unity must satisfy $P(z_0)=0$. But $(x^5-1)/(x-1)$ has exactly all 5th roots of unity excluding 1, and $(x^5-1)/(x-1)=x^4+x^3+x^2+x+1$. Thus this must divide $P$ which implies $P(x)=4(x^4+x^3+x^2+x+1)$. This completes the proof.

 



\textbf{Solution 2}

First, assume that $z_0\in \mathbb{R}$, so $z_0=1$ or $-1$. $1$ does not work because $P(1)\geq 4$. Assume that $z_0=-1$. Then $0=P(-1)=4-a+b-c+d$,  we have $4+b+d=a+c\leq 4+b$, so $d=0$. Also, $a=4$ has to be true since $4+b=a+c \leq a+b$. Now $4+b=4+c$ gives $b=c$, therefore the only possible choices for $(a,b,c,d)$ are $(4,t,t,0)$. In these cases, $P(-1)=4-4+t-t+0=0$. The sum of $P(1)$ over these cases is $\sum_{t=0}^{4} (4+4+t+t) = 40+20=60$.

 Second, assume that $z_0\in \mathbb{C} \backslash \mathbb{R}$, so $z_0=x_0+iy_0$ for some real $x_0, y_0$, $|x_0|<1$. By conjugate roots theorem we have that $P(z_0)=P(z_0^{*})=0$, therefore $(z-z_0)(z-z_0^{*}) = (z^2 - 2x_0*z + 1)$ is a factor of $P(z)$, and we may assume that 

 
 \[P(z) = (z^2-2x_0 z + 1)(4z^2 + pz + d)\]

 for some real $p$. Expanding this polynomial and comparing the coefficients, we have the following equations:

 
 \[p-8x_0 = a\]
 \[d+4-2px_0 = b\]
 \[p-2dx_0 = c\]

 From the first and the third we may deduce that $2x_0 = \frac{a-c}{d-4}$ and that $p=\frac{da-4c}{d-4}$, if $d\neq 4$ (we will consider $d=4$ by the end). Let $k=2px_0=\frac{(a-c)(da-4c)}{(4-d)^2}$. From the second equation, we know that $k=d+4-b$ is non-negative.

 Consider the following cases:

 Case 1: $a=c$. Then $k=0$, $b=d+4$, so $a=b=c=4$, $d=0$. However, this has already been found (i.e. the form of $(4,t,t,0)$).

 Case 2: $a>c\geq 0$. Then since $k\geq 0$, we have $da-4c\geq 0$. However, $da \leq 4c$, therefore $da-4c=0$. This is true only when $d=c$. Also, we get $k=0$ again. In this case, $b=d+4$, so $a=b=4$, $c=d=0$, $x_0=-1/2$. $P(z)$ has a root $z_0=e^{i2\pi/3}$. $P(1)=12$.

 Last case: $d=4$. We have $a=b=c=d=4$ and that $P(z)$ has a root $z_0=e^{i2\pi/5}$. $P(1)=20$.

 Therefore the desired sum is $60+12+20=92 ...\framebox{B}$.

 



\textbf{Solution 3}

First, notice that $z=1$ cannot be a root of the polynomial because $a,b,c,d \geq 0$. Multiplying the polynomial by $(z-1)$ yields $P(z)(z-1) = 4z^5-(4-a)z^4-(a-b)z^3-(b-c)z^2-(c-d)z-d$, so for $z \neq 1$ to be a root of $P(z)$, $4z^5 = (4-a)z^4+(a-b)z^3+(b-c)z^2+(c-d)z+d$. Now we consider the root $z_0$ with $|z_0|=1$. $|4z_0^5| = 4$, so the right hand side must have absolute value 4. By the triangle inequality, $|(4-a)z_0^4+(a-b)z_0^3+(b-c)z_0^2+(c-d)z_0+d|$ $\leq |(4-a)z_0^4|+|(a-b)z_0^3| + |(b-c)z_0^2| + |(c-d)z_0|+d$ $= (4-a)+(a-b)+(b-c)+(c-d)+d=4$, with equality if and only if each of $4z_0^5$, $(4-a)z_0^4$, $(a-b)z_0^3$, $(b-c)z_0^2$, $(c-d)z_0$, and $d$ is either zero or in the same direction as all the others when looked at as vectors in the complex plane. 

 We can now divide into two cases: $d \neq 0$ and $d=0$. If $d \neq 0$, then $4z_0^5$ must be real by the previous argument, so $z_0$ is a fifth root of unity. Also, $(4-a)z_0^4$, $(a-b)z_0^3$, $(b-c)z_0^2$, and $(c-d)z_0$ must all be zero because if $z_0$ is a fifth root of unity, none of these can be real numbers with positive absolute value. Therefore, $a=4$, $b=a$, $c=b$, and $d=c$, leading to the solution $(a,b,c,d)=(4,4,4,4)$. Just to be sure, we can easily verify that this solution leads to the six complex numbers under question being in the same direction.

 If $d=0$, then each of $4z_0^5$, $(4-a)z_0^4$, $(a-b)z_0^3$, $(b-c)z_0^2$, and $cz_0$ must either be zero or in the same direction as all the others, so each of $4z_0^4$, $(4-a)z_0^3$, $(a-b)z_0^2$, $(b-c)z_0$, and $c$ must either be zero or in the same direction as all the others. We can divide this into two cases: $c \neq 0$ and $c=0$. If $c \neq 0$, then $4z_0^4$ must be real. Then, $z_0$ is a fourth root of unity. If $z_0$ is not a second root of unity, $(4-a)z_0^3$, $(a-b)z_0^2$, and $(b-c)z_0$ must be zero, implying that $a=4$, $b=a=4$, and $c=b=4$, leading to the solution $(a,b,c,d)=(4,4,4,0)$. If $z_0$ is also a second root of unity, $(4-a)z_0^3$ and $(b-c)z_0$ must be zero but $(a-b)z_0^2$ can be anything. This implies $a=4$ and $b=c$ with no other restrictions, leading to the new solutions $(a,b,c,d) = (4,3,3,0), (4,2,2,0), (4,1,1,0), (4,0,0,0)$. 

 If $c=0$, then we can similarly show that each of $4z_0^3$, $(4-a)z_0^2$, $(a-b)z_0$, and $b$ must be zero or in the same direction as all the others. If $b \neq 0$, then $z_0$ must be a third root of unity, so $(4-a)z_0^2$ and $(a-b)z_0$ must be zero, implying $a=b=4$, leading to the new solution $(a,b,c,d)=(4,4,0,0)$. 

 If $b=0$, then we can similarly show that each of $4z_0^2$, $(4-a)z_0$, and $a$ must be zero or in the same direction as the others. For $|z_0|=1$, $a = 4$, but we have already counted the solution $(a,b,c,d)=(4,0,0,0)$. 

 Then, the complete list of solutions is $(a,b,c,d)=(4,4,4,4),(4,4,4,0),(4,3,3,0),(4,2,2,0),(4,1,1,0),(4,0,0,0),(4,4,0,0)$, leading to a sum of $\framebox{B}=92$.

 


	\item (\textit{AMC 12 12B 2012}) Define the function $f_1$ on the positive integers by setting $f_1(1)=1$ and if $n=p_1^{e_1}p_2^{e_2}\cdots p_k^{e_k}$ is the prime factorization of $n>1$, then 
 \[f_1(n)=(p_1+1)^{e_1-1}(p_2+1)^{e_2-1}\cdots (p_k+1)^{e_k-1}.\]
For every $m\ge 2$, let $f_m(n)=f_1(f_{m-1}(n))$. For how many $N$s in the range $1\le N\le 400$ is the sequence $(f_1(N),f_2(N),f_3(N),\dots )$ unbounded?


 \textbf{Note:} A sequence of positive numbers is unbounded if for every integer $B$, there is a member of the sequence greater than $B$.


 $\textbf{(A)}\ 15\qquad\textbf{(B)}\ 16\qquad\textbf{(C)}\ 17\qquad\textbf{(D)}\ 18\qquad\textbf{(E)}\ 19$

 \textbf{Answer: }D



\textbf{Solution 1}

First of all, notice that for any odd prime $p$, the largest prime that divides $p+1$ is no larger than $\frac{p+1}{2}$, therefore eventually the factorization of $f_k(N)$ does not contain any prime larger than $3$. Also, note that $f_2(2^m) = f_1(3^{m-1})=2^{2m-4}$, when $m=4$ it stays the same but when $m\geq 5$ it grows indefinitely. Therefore any number $N$ that is divisible by $2^5$ or any number $N$ such that $f_k(N)$ is divisible by $2^5$ makes the sequence $(f_1(N),f_2(N),f_3(N),\dots )$ unbounded. There are $12$ multiples of $2^5$ within $400$. $2^4 5^2=400$ also works: $f_2(2^4 5^2) = f_1(3^4 \cdot 2) = 2^6$.

 Now let's look at the other cases. Any first power of prime in a prime factorization will not contribute the unboundedness because $f_1(p^1)=(p+1)^0=1$. At least a square of prime is to contribute. So we test primes that are less than $\sqrt{400}=20$:

 $f_1(3^4)=4^3=2^6$ works, therefore any number $\leq 400$ that are divisible by $3^4$ works: there are $4$ of them.

 $3^3 \cdot Q^2$ could also work if $Q^2$ satisfies $2~|~f_1(Q^2)$, but $3^3 \cdot 5^2 > 400$.

 $f_1(5^3)=6^2 = 2^2 3^2$ does not work.

 $f_1(7^3)=8^2=2^6$ works. There are no other multiples of $7^3$ within $400$.

 $7^2 \cdot Q^2$ could also work if $4~|~f_1(Q^2)$, but $7^2 \cdot 3^2 > 400$ already.

 For number that are only divisible by $p=11, 13, 17, 19$, they don't work because none of these primes are such that $p+1$ could be a multiple of $2^5$ nor a multiple of $3^4$.

 In conclusion, there are $12+1+4+1=18$ number of $N$'s ... $\framebox{D}$.

 



\textbf{Solution 2}

Say a number $n$ is $\textit{boring}$ if the sequence $\{f_k(n)\}$ is bounded; otherwise $n$ is $\textit{interesting}$. 

 It is clear that $n$ is interesting iff $f_1(n)$ is. For a prime, $p$, $f_1(p)=1$. Thus, for a prime, $p_k$, in the prime factorization of $n$; if $e_k=1$ then $f_1(n)=f_1(n/p_k)$, so $n$ is interesting iff $n/p_k$ is. Continuing in this manner, we can divide $n$ by all such primes $p_k$ for which $e_k=1$; and $n$ is interesting iff each of these resulting numbers are. Finally we will end up with a number  whose prime factorization  contains only exponents $\ge 2$. Let $S$ be the set of such numbers. It suffices to find all interesting numbers in $S$; all other interesting numbers will be multiples of these. Note that  $f_1(n)\mid f_1(kn)$ for all $k\in\mathbb{N}$; so by induction a multiple of an interesting number is also interesting.

 The set $S$ contains either powers of primes $\le 19$, or a product of two such powers.

 We have $f_2(2^a)=f_1(3^{a-1}) = 2^{2(a-2)}$; since $2(a-2)> a$ iff $a\ge 5$, so $2^5$ and its multiples are interesting; there are $12$ such.

 We have $f_2(3^b)=f_1(2^{2(b-1)}) = 3^{2b-3}$; since $2b-3> b$ iff $b\ge 4$, so $3^4$ and its multiples are interesting; there are $4$ such.

 Among the remaining prime powers only $7^3 \textbf{ is interesting}$ but it has no other multiples in $[400]$.

 We are now left with numbers $n\in S$ that are products of two prime powers, i.e $n=p^aq^b$, $a,b \ge 2$.  

 We have $f_2(2^a3^b)= 2^{2(a-2)}3^{2b-3}$, so a number of this form is interesting iff $a\ge 5$ or $b\ge 4$; they are already counted above.

 If $5^2\mid n$ then $n=2^2\cdot 5^2$ or $2^3\cdot 5^2$ or $3^2\cdot 5^2$ (all boring) or $n=2^4\cdot 5^2=400$, $\textbf{which is interesting}$, but has no other multiples. 

 If $7^2\mid n$ then $n$ is $2^2\cdot 7^2$ or $2^3\cdot 7^2$ and both are boring. 

 No larger prime can divide $n$ because even $11^2\cdot 2^2 > 400$. 

 We have found all the interesting numbers: $12$ multiples of $2^5$, $4$ multiples of $3^4$, $2^4\cdot 5^2$, and $7^3$ for a total of $12+4+1+1= 18$ ... $\framebox{D}$.

 


	\item (\textit{AMC 12 12B 2012}) Let $S=\{(x,y) : x\in \{0,1,2,3,4\}, y\in \{0,1,2,3,4,5\},\text{ and } (x,y)\ne (0,0)\}$. 
Let $T$ be the set of all right triangles whose vertices are in $S$. For every right triangle $t=\triangle{ABC}$ with vertices $A$, $B$, and $C$ in counter-clockwise order and right angle at $A$, let $f(t)=\tan(\angle{CBA})$. What is 
 \[\prod_{t\in T} f(t)?\]
$\textbf{(A)}\ 1\qquad\textbf{(B)}\ \frac{625}{144}\qquad\textbf{(C)}\ \frac{125}{24}\qquad\textbf{(D)}\ 6\qquad\textbf{(E)}\ \frac{625}{24}$

 \textbf{Answer: }B



\textbf{Solution 1}

Consider reflections. For any right triangle $ABC$ with the right labeling described in the problem, any reflection $A'B'C'$ labeled that way will give us $\tan CBA \cdot \tan C'B'A' = 1$. First we consider the reflection about the line $y=2.5$. Only those triangles $\subseteq T$ that have one vertex at $(0,5)$ do not reflect to a traingle $\subseteq T$. Within those triangles, consider a reflection about the line $y=5-x$. Then only those triangles $\subseteq T$ that have one vertex on the line $y=0$ do not reflect to a triangle $\subseteq T$. So we only need to look at right triangles that have vertices $(0,5), (*,0), (*,*)$. There are three cases:

 Case 1: $A=(0,5)$. Then $B=(*,0)$ is impossible.

 Case 2: $B=(0,5)$. Then we look for $A=(x,y)$ such that $\angle BAC=90^{\circ}$ and that $C=(*,0)$. They are: $(A=(x,5), C=(x,0))$,  $(A=(3,2), C=(1,0))$ and $(A=(4,1), C=(3,0))$. The product of their values of $\tan \angle CBA$ is $\frac{5}{1}\cdot  \frac{5}{2} \cdot \frac{5}{3} \cdot \frac{5}{4} \cdot \frac{1}{4} \cdot \frac{2}{3} = \frac{625}{144}$.

 Case 3: $C=(0,5)$. Then $A=(*,0)$ is impossible.

 Therefore $\boxed{\textbf{(B)}  \ \frac{625}{144}}$ is the answer.

 



\textbf{Solution 2}

This is just another way for the reasoning of solution 1. Picture the question to be a grid of unit squares instead of a coordinate system. Note that the restriction, (This is just to make visualization easier.) Define a "cell" to be a rectangle in the set of $S.$ For example, a cell can be (labeled in a red):

 
                    \begin{figure}[H]
                        \centering
                        \includegraphics[width=0.5\textwidth]{Solution/AMC_12_12B_2012_25_Solution_Figure_1.png}
                    \end{figure}
                    

 Note that choosing any three points and labeling them according to the problem will result in a product of one. For example, with the cell we just labeled, the four triangles we can create are:

 
                    \begin{figure}[H]
                        \centering
                        \includegraphics[width=0.5\textwidth]{Solution/AMC_12_12B_2012_25_Solution_Figure_2.png}
                    \end{figure}
                    

                    \begin{figure}[H]
                        \centering
                        \includegraphics[width=0.5\textwidth]{Solution/AMC_12_12B_2012_25_Solution_Figure_3.png}
                    \end{figure}
                    

                    \begin{figure}[H]
                        \centering
                        \includegraphics[width=0.5\textwidth]{Solution/AMC_12_12B_2012_25_Solution_Figure_4.png}
                    \end{figure}
                    

                    \begin{figure}[H]
                        \centering
                        \includegraphics[width=0.5\textwidth]{Solution/AMC_12_12B_2012_25_Solution_Figure_5.png}
                    \end{figure}
                    

 If we define the longer side to be $x$ and the shorter side to be $y,$ then the product will be $\frac{x}{y} \cdot \frac{y}{x} \cdot \frac{x}{y} \cdot \frac{y}{x}=1,$ and we are done.

 Otherwise, the three points are not contained in a "cell." This will result in the solution 1 path as described before. Our three points must take the form $(0,5), (*,0), (*,*),$ where $*$ is a number defined by the boundaries of $S.$ Thus, by the three cases, our answer is $\boxed{\textbf{(B)}  \ \frac{625}{144}}.$



\end{enumerate}